% Draft retention material only; not automatically included.
% Integrate into the appropriate sections of paper/main.tex.
% Current notation uses integer N,h,Y and C^*=0 for h>=N.
% No arbitrary-translate spectral large-sieve input is asserted.

% Place after the square/endpoint-duality lemma.
\begin{proposition}[A terminal-range version of the collar implication]
\label{prop:terminal-collar-retained}
Fix \(\eta>0\). Suppose that for every \(\varepsilon>0\),
\[
 \mathcal E_\lambda(N,H)
 \ll_{\varepsilon,\eta}\frac{N^{2+\varepsilon}}H
 \qquad
 (H\text{ dyadic},\ 1\le H\le N^{1-\eta}),
\]
and
\[
 \sum_{N^{1-\eta}<h<N}C^*_{\lambda,h}(N)
       \ll_{\varepsilon,\eta}N^{1+\varepsilon}.
\]
Then, for every \(\varepsilon>0\),
\[
 \sup_{\alpha\in\mathbb R}
 \left|\sum_{n\le N}\lambda(n)\mathrm e(n\alpha)\right|
       \ll_{\varepsilon,\eta}N^{1/2+\varepsilon}.
\]
In particular the summatory Liouville criterion gives RH.
\end{proposition}
\begin{proof}
The diagonal in the square identity and the shift \(h=1\) contribute
\(O(N)\). For each dyadic \(H\le N^{1-\eta}\),
\[
 \sum_{H<h\le2H}C^*_{\lambda,h}(N)
 \le\sqrt H\,\mathcal E_\lambda(N,H)^{1/2}
 \ll_{\varepsilon,\eta}N^{1+\varepsilon/2}.
\]
There are \(O(\log(2N))\) such blocks. Include the whole block that
crosses \(N^{1-\eta}\), if one is present. Every shift \(h\ge2\) not
covered is greater than \(N^{1-\eta}\); its contribution is bounded
by the terminal-range hypothesis. Thus
\[
 |A_N(\alpha)|^2
 \ll_{\varepsilon,\eta}
 N+N^{1+\varepsilon/2}\log(2N)+N^{1+\varepsilon}.
\]
Choose a smaller input exponent and take a square root.
\end{proof}

% Place immediately after the exact coefficient-uniform maximum.
\begin{corollary}[The original explicit coefficient-uniform gap]
\label{cor:original-coefficient-gap}
Fix \(\eta\in(0,3/4)\). For each sufficiently large integer \(N\), choose
a dyadic \(H\) with
\[
 \frac12N^{1/4}<H\le N^{1/4}\le N^{1-\eta}.
\]
The constant sequence \(a_n=1\) satisfies
\[
 \mathcal E_{\mathbf1}(N,H)\ge\frac1{16}N^{9/4}.
\]
A coefficient-uniform collar estimate with \(\varepsilon=1/8\)
would instead bound this by
\[
 C_{1/8}N^{17/8}H^{-1}\le2C_{1/8}N^{15/8}.
\]
The ratio to the proposed scale is at least \(N^{3/8}/32\);
no fixed coefficient-uniform constant can work.
\end{corollary}
\begin{proof}
For large \(N\), \(2H\le N/2\), so
\[
 \mathcal E_{\mathbf1}(N,H)
 =\sum_{H<h\le2H}(N-h)^2
 \ge\frac{HN^2}4>\frac18N^{9/4}\ge\frac1{16}N^{9/4}.
\]
The lower bound on \(H\) gives the upper bound for the proposed scale.
Finally \(9/4-15/8=3/8\).
\end{proof}

% Restore the exponent figure separately with fixed 0<u<1 and H~N^u.
% Curves: worst-case 2+u and zero-loss target 2-u.
% Optional dashed line 2+1/8-u distinguishes the explicit corollary.
% At u=1 do not imply that the empty block H=N has energy N^3.

% Place beside the original correlation discussion.
\begin{lemma}[Square-divisor expansion of a correlation]
\label{lem:square-divisor-correlation-retained}
For every positive integer \(n\),
\[
 \lambda(n)=\sum_{r^2\mid n}\mu(n/r^2).
\]
Consequently, for integers \(Y\ge0\) and \(h\ge1\),
\begin{align*}
 C_{\lambda,h}(Y)
 &=\sum_{r\le\sqrt Y}\sum_{s\le\sqrt{Y+h}}
   \sum_{\substack{1\le n\le Y\\r^2\mid n,\ s^2\mid n+h}}
       \mu(n/r^2)\mu((n+h)/s^2)\\
 &=\sum_{\substack{r\le\sqrt Y,\ s\le\sqrt{Y+h}\\
                   \gcd(r,s)^2\mid h}}
   \ \sum_{\substack{m,k\ge1\\r^2m\le Y\\s^2k-r^2m=h}}
       \mu(m)\mu(k).
\end{align*}
The \((r,s)=(1,1)\) stratum is
\(\sum_{n\le Y}\mu(n)\mu(n+h)\).
\end{lemma}
\begin{proof}
Both sides of the first identity are multiplicative. At \(n=p^e\),
only \(e-2j=0\) for even \(e\), or \(e-2j=1\) for odd \(e\),
contributes. The values are respectively \(1\) and \(-1\).
Insert the identity in each factor of the finite correlation and
interchange finite sums. Put \(n=r^2m\) and \(n+h=s^2k\).
The congruences \(n=0\bmod r^2\), \(n=-h\bmod s^2\) are compatible
exactly when \(\gcd(r^2,s^2)=\gcd(r,s)^2\) divides \(h\).
When compatible, they have one residue class modulo
\(\operatorname{lcm}(r^2,s^2)\); its displayed finite range may be empty.
The leading stratum is immediate.
\end{proof}

The expansion retains the signed arithmetic data. Other strata cannot
be bounded by the leading one: they involve two affine forms with
parameters depending on \(r,s,h\). The coefficient-uniform construction
excludes discarding this information and using only \(|\mu|\le1\).

% Missing original bibliography entries (19).
% Insert inside the existing thebibliography, not a second environment.
% Six original sources already survive under keys:
% titchmarsh, di, mr-short, mrt, tao, pt.

\bibitem{apostol}
T. M. Apostol, \emph{Introduction to Analytic Number Theory},
Undergraduate Texts in Mathematics, Springer-Verlag, New York, 1976.

\bibitem{iwaniec-kowalski}
H. Iwaniec and E. Kowalski, \emph{Analytic Number Theory},
American Mathematical Society Colloquium Publications \textbf{53},
American Mathematical Society, Providence, RI, 2004.

\bibitem{iwaniec-spectral}
H. Iwaniec, \emph{Spectral Methods of Automorphic Forms}, 2nd ed.,
Graduate Studies in Mathematics \textbf{53},
American Mathematical Society/Revista Matem\'atica Iberoamericana, 2002.

\bibitem{friedlander-iwaniec}
J. B. Friedlander and H. Iwaniec, \emph{Opera de Cribro},
American Mathematical Society Colloquium Publications \textbf{57},
American Mathematical Society, Providence, RI, 2010.

\bibitem{bombieri}
E. Bombieri, \emph{Le grand crible dans la th\'eorie analytique des nombres},
Ast\'erisque \textbf{18}, Soci\'et\'e Math\'ematique de France,
2nd ed., 1987.

\bibitem{montgomery}
H. L. Montgomery, \emph{Ten Lectures on the Interface Between Analytic
Number Theory and Harmonic Analysis},
CBMS Regional Conference Series in Mathematics \textbf{84},
American Mathematical Society, 1994.

\bibitem{selberg}
A. Selberg, \emph{On the zeros of Riemann's zeta-function},
Skr. Norske Vid. Akad. Oslo I \textbf{10} (1942), 1--59.

\bibitem{odlyzko-teriele}
A. M. Odlyzko and H. J. J. te Riele,
\emph{Disproof of the Mertens conjecture},
J. Reine Angew. Math. \textbf{357} (1985), 138--160.

\bibitem{sarnak-letter}
P. Sarnak, \emph{Three Lectures on the M\"obius Function, Randomness
and Dynamics}, lecture notes, Institute for Advanced Study, Princeton, 2011.

\bibitem{bourgain-sarnak-ziegler}
J. Bourgain, P. Sarnak and T. Ziegler, \emph{Disjointness of M\"obius
from horocycle flows}, in \emph{From Fourier Analysis and Number Theory
to Radon Transforms and Geometry}, Developments in Mathematics
\textbf{28}, Springer, New York, 2013, pp. 67--83.

\bibitem{matomaki-radziwill-tao-signs}
K. Matom\"aki, M. Radziwi\l\l\ and T. Tao,
\emph{Sign patterns of the M\"obius and Liouville functions},
Q. J. Math. \textbf{67} (2016), 461--482.

\bibitem{tao-teravainen}
T. Tao and J. Ter\"av\"ainen, \emph{The structure of logarithmically
averaged correlations of multiplicative functions, with applications
to the Chowla and Elliott conjectures},
Duke Math. J. \textbf{168} (2019), 1977--2027.

\bibitem{klurman}
O. Klurman, \emph{Correlations of multiplicative functions and
applications}, Compos. Math. \textbf{153} (2017), 1622--1657.

\bibitem{teravainen-binary}
J. Ter\"av\"ainen, \emph{On binary correlations of multiplicative
functions}, Forum Math. Sigma \textbf{9} (2021), e16.

\bibitem{teravainen-walker}
J. Ter\"av\"ainen and A. Walker, \emph{On a Bohr set analogue of
Chowla's conjecture}, Math. Z. \textbf{310} (2025), article 75,
\href{https://doi.org/10.1007/s00209-025-03770-2}
 {doi:10.1007/s00209-025-03770-2}.

\bibitem{chavez}
G. Chavez, \emph{Correlations of the M\"obius and Liouville functions
with their partial sums}, Int. J. Number Theory (2026),
\href{https://doi.org/10.1142/S1793042126500934}
 {doi:10.1142/S1793042126500934}.
Preprint under the title \emph{Correlations of multiplicative functions
with their partial sums},
\href{https://arxiv.org/abs/2409.02106v10}{arXiv:2409.02106v10}.

\bibitem{krishnamoorthy-chowla}
K. Krishnamoorthy, \emph{On variants of Chowla's conjecture},
\href{https://arxiv.org/abs/2501.10962v2}{arXiv:2501.10962v2} (2025).
Published version:
\href{https://doi.org/10.1090/proc/17544}{doi:10.1090/proc/17544}.
% Primary arXiv metadata confirms the related DOI.
% Publisher volume/page metadata was inaccessible in this audit.

\bibitem{pascadi}
A. Pascadi, \emph{Large sieve inequalities for exceptional Maass forms
and the greatest prime factor of \(n^2+1\)},
Forum Math. Pi \textbf{14} (2026), e8, 1--54,
\href{https://doi.org/10.1017/fmp.2026.10025}
 {doi:10.1017/fmp.2026.10025}.

\bibitem{granville-soundararajan}
A. Granville and K. Soundararajan, \emph{Large character sums:
pretentious characters and the P\'olya--Vinogradov theorem},
J. Amer. Math. Soc. \textbf{20} (2007), 357--384.
